% !TeX TXS-program:compile = txs:///arara
% arara: pdflatex: {shell: no, synctex: no, interaction: batchmode}
% arara: pdflatex: {shell: no, synctex: no, interaction: batchmode}

\documentclass{article}
\usepackage[margin=1in]{geometry}
\usepackage{vscodeicons}
\usepackage{longtable}
\usepackage{booktabs}
\usepackage{hyperref}
\usepackage{fancyvrb}
\usepackage{shortvrb}
\MakeShortVerb{\|}
\setlength{\parindent}{0pt}

\newenvironment{vscodeiconscase}
{%
	\begin{longtable}{cllr}%icon + name + macro + page
		\cmidrule[\heavyrulewidth]{1-4}% \toprule
		\bfseries Icon & \bfseries Name & \bfseries Direct command & \bfseries PDF page \\
		\cmidrule{1-4}
		\endhead
}%
{%
		\cmidrule[\heavyrulewidth]{1-4}
	\end{longtable}%
}

\NewDocumentCommand{\showcasevscodeicondefault}{mm}{%name + macro + page
	\Large\vscodeicon{#1} & \itshape #1 & \ttfamily \textbackslash vscodeicon\{#1\} & \sffamily #2 \\
}

\NewDocumentCommand{\showcasevscodeiconfiletype}{mm}{%name + macro + page
	\Large\vscodeicon[filetype]{#1} & \itshape #1 & \ttfamily \textbackslash vscodeicon[filetype]\{#1\} & \sffamily #2 \\
}

\NewDocumentCommand{\showcasevscodeiconfoldertype}{mm}{%name + macro + page
	\Large\vscodeicon[foldertype]{#1} & \itshape #1 & \ttfamily \textbackslash vscodeicon[foldertype]\{#1\} & \sffamily #2 \\
}

\begin{document}

\title{The \textsf{vscodeicons} package (0.1.0) -- \today}
\author{%
	Cédric Pierquet\\%
	\url{https://github.com/vscode-icons/vscode-icons} (Source version v12.15.0)\\%
	\url{https://github.com/cpierquet/latex-packages/tree/main/vscodeicons} (\LaTeX{} package)
}
\date{cpierquet -- at -- outlook . fr}\maketitle

This package provides \LaTeX{} support for the vscode icons project (\url{https://github.com/vscode-icons/vscode-icons}).

%\medskip
%
%\hspace*{5mm}\fbox{\begin{minipage}{0.75\linewidth}{\emph{A set of 5944 free MIT-licensed high-quality SVG/PDF icons for you to use in your web projects. Each icon is designed on a 24x24 grid and a 2px stroke.}}\end{minipage}}

\medskip

To use \textsf{vscodeicons} in your document, include the package with |\usepackage{vscodeicons}|.

An icon can be accessed using the icon name. A list of all included icons with their respective \textit{page} can be found at the end of this document.

\section{Example}

\begin{Verbatim}[frame=single]
...
\usepackage{vscodeicons}
...
\begin{document}
...
Inline\vscodeicon{folder}version
...
\end{document}
\end{Verbatim}

Result: Inline\vscodeicon{folder}version

\section{Usage}

\subsection{Generic macro}

\begin{Verbatim}[frame=single]
\vscodeicon%
  [filetype=TF,foldertype=TF,height=...,(d)strut=...]%
  <includegraphics options>%
  {name}
\end{Verbatim}

The boolean \texttt{[filetype]} activate \textsf{filetype} version of icon (if available) and the boolean \texttt{[foldertype]} activate \textsf{foldertype} version of icon (if available).

And key \texttt{[height=...]} can be:

\begin{itemize}
	\item \texttt{dauto} (by default): inline insertion with height/depth;
	\item \texttt{auto}: inline insertion with height/nodepth;
	\item \texttt{a length}: normal insertion with fixed height;
	\item \texttt{manuel height/depth}: normal insertion with fixed height/depth.
\end{itemize}

The keys \texttt{[(d)strut=...]} can be used to set the \textit{reference} box for dimension calculations.

\medskip

By default, icons are given with a small border, so the given height is the height of the \textit{full} box, not the hight of the icons.

\medskip

Example: Inline{\setlength\fboxsep{0pt}\fbox{\vscodeicon{root-folder}}}version or inline {\setlength\fboxsep{0pt}\fbox{\vscodeicon{root-folder}}} version.

\pagebreak

\subsection{Examples}

\begin{Verbatim}[frame=single]
%normal version: inline with automatic height/depth (from qX characters)
\vscodeicon[filetype,height=dauto]{arduino} %or \vscodeicon[filetype]{arduino}
\end{Verbatim}

{\Huge q\vscodeicon[filetype]{arduino}X}

\begin{Verbatim}[frame=single]
%normal version: inline with automatic height/depth (from specified characters)
(q\vscodeicon[foldertype,height=dauto,dstrut=(qB)]{story}B)
\end{Verbatim}

{\Huge (q\vscodeicon[foldertype,height=dauto,dstrut=(qB)]{story}B)}

\begin{Verbatim}[frame=single]
%normal version: inline with automatic height (from X character) default version
\vscodeicon[foldertype,height=auto]{story}
\end{Verbatim}

{\Huge X\vscodeicon[foldertype,height=auto]{story}X}

\begin{Verbatim}[frame=single]
%normal version: inline with manual height/depth
\vscodeicon[foldertype,height={0.8em/-0.4em}]{story}
\end{Verbatim}

{\Huge qb\vscodeicon[foldertype,height={0.8em/-0.4em}]{story}Xz}

\begin{Verbatim}[frame=single]
%fixed height
\vscodeicon[height=2in]{folder-opened}
\end{Verbatim}

\vscodeicon[height=2in]{folder-opened}

\subsection{Used packages}

\texttt{ifthen}, \texttt{calc}, \texttt{graphicx}, \texttt{xstring} and \texttt{simplekv} are loaded and used by the package.

\subsection{Bugs}

For bug reports and feature requests, report on the GitHub repository \url{https://github.com/cpierquet/latex-packages/issues}.

\newpage

\section{List of icons}

\subsection{Default version (5 icons)}

The \textsf{PDF} file is \texttt{vscodeicons-default-all.pdf}

\begin{vscodeiconscase}
\showcasevscodeicondefault{file}{1}
\showcasevscodeicondefault{folder-opened}{2}
\showcasevscodeicondefault{folder}{3}
\showcasevscodeicondefault{root-folder-opened}{4}
\showcasevscodeicondefault{root-folder}{5}
\end{vscodeiconscase}

\subsection{FileType version (1061 icons)}

The \textsf{PDF} file is \texttt{vscodeicons-filetype-all.pdf}

\begin{vscodeiconscase}
\showcasevscodeiconfiletype{access}{1}
\showcasevscodeiconfiletype{access2}{2}
\showcasevscodeiconfiletype{actionscript}{3}
\showcasevscodeiconfiletype{actionscript2}{4}
\showcasevscodeiconfiletype{ada}{5}
\showcasevscodeiconfiletype{advpl}{6}
\showcasevscodeiconfiletype{affectscript}{7}
\showcasevscodeiconfiletype{affinitydesigner}{8}
\showcasevscodeiconfiletype{affinityphoto}{9}
\showcasevscodeiconfiletype{affinitypublisher}{10}
\showcasevscodeiconfiletype{agda}{11}
\showcasevscodeiconfiletype{agents}{12}
\showcasevscodeiconfiletype{ai}{13}
\showcasevscodeiconfiletype{ai2}{14}
\showcasevscodeiconfiletype{al}{15}
\showcasevscodeiconfiletype{allcontributors}{16}
\showcasevscodeiconfiletype{alloy}{17}
\showcasevscodeiconfiletype{al-dal}{18}
\showcasevscodeiconfiletype{angular}{19}
\showcasevscodeiconfiletype{ansible}{20}
\showcasevscodeiconfiletype{antlers-html}{21}
\showcasevscodeiconfiletype{antlr}{22}
\showcasevscodeiconfiletype{anyscript}{23}
\showcasevscodeiconfiletype{apache}{24}
\showcasevscodeiconfiletype{apex}{25}
\showcasevscodeiconfiletype{apib}{26}
\showcasevscodeiconfiletype{apib2}{27}
\showcasevscodeiconfiletype{api-extractor}{28}
\showcasevscodeiconfiletype{apl}{29}
\showcasevscodeiconfiletype{applescript}{30}
\showcasevscodeiconfiletype{appscript}{31}
\showcasevscodeiconfiletype{appsemble}{32}
\showcasevscodeiconfiletype{appveyor}{33}
\showcasevscodeiconfiletype{arduino}{34}
\showcasevscodeiconfiletype{asciidoc}{35}
\showcasevscodeiconfiletype{aseprite}{36}
\showcasevscodeiconfiletype{asp}{37}
\showcasevscodeiconfiletype{aspx}{38}
\showcasevscodeiconfiletype{assembly}{39}
\showcasevscodeiconfiletype{astro}{40}
\showcasevscodeiconfiletype{astroconfig}{41}
\showcasevscodeiconfiletype{atom}{42}
\showcasevscodeiconfiletype{ats}{43}
\showcasevscodeiconfiletype{attw}{44}
\showcasevscodeiconfiletype{audio}{45}
\showcasevscodeiconfiletype{aurelia}{46}
\showcasevscodeiconfiletype{autohotkey}{47}
\showcasevscodeiconfiletype{autoit}{48}
\showcasevscodeiconfiletype{avif}{49}
\showcasevscodeiconfiletype{avro}{50}
\showcasevscodeiconfiletype{awk}{51}
\showcasevscodeiconfiletype{aws}{52}
\showcasevscodeiconfiletype{azure}{53}
\showcasevscodeiconfiletype{azurepipelines}{54}
\showcasevscodeiconfiletype{babel}{55}
\showcasevscodeiconfiletype{babel2}{56}
\showcasevscodeiconfiletype{bak}{57}
\showcasevscodeiconfiletype{ballerina}{58}
\showcasevscodeiconfiletype{bat}{59}
\showcasevscodeiconfiletype{bats}{60}
\showcasevscodeiconfiletype{bazaar}{61}
\showcasevscodeiconfiletype{bazel}{62}
\showcasevscodeiconfiletype{bazel-ignore}{63}
\showcasevscodeiconfiletype{bazel-version}{64}
\showcasevscodeiconfiletype{befunge}{65}
\showcasevscodeiconfiletype{bicep}{66}
\showcasevscodeiconfiletype{biml}{67}
\showcasevscodeiconfiletype{binary}{68}
\showcasevscodeiconfiletype{biome}{69}
\showcasevscodeiconfiletype{bitbucketpipeline}{70}
\showcasevscodeiconfiletype{bithound}{71}
\showcasevscodeiconfiletype{blade}{72}
\showcasevscodeiconfiletype{blender}{73}
\showcasevscodeiconfiletype{blitzbasic}{74}
\showcasevscodeiconfiletype{bolt}{75}
\showcasevscodeiconfiletype{bosque}{76}
\showcasevscodeiconfiletype{bower}{77}
\showcasevscodeiconfiletype{bower2}{78}
\showcasevscodeiconfiletype{browserslist}{79}
\showcasevscodeiconfiletype{bruno}{80}
\showcasevscodeiconfiletype{buckbuild}{81}
\showcasevscodeiconfiletype{buf}{82}
\showcasevscodeiconfiletype{bun}{83}
\showcasevscodeiconfiletype{bundlemon}{84}
\showcasevscodeiconfiletype{bundler}{85}
\showcasevscodeiconfiletype{bunfig}{86}
\showcasevscodeiconfiletype{c}{87}
\showcasevscodeiconfiletype{c2}{88}
\showcasevscodeiconfiletype{c3}{89}
\showcasevscodeiconfiletype{cabal}{90}
\showcasevscodeiconfiletype{caddy}{91}
\showcasevscodeiconfiletype{cake}{92}
\showcasevscodeiconfiletype{cakephp}{93}
\showcasevscodeiconfiletype{capacitor}{94}
\showcasevscodeiconfiletype{capnp}{95}
\showcasevscodeiconfiletype{cargo}{96}
\showcasevscodeiconfiletype{casc}{97}
\showcasevscodeiconfiletype{cddl}{98}
\showcasevscodeiconfiletype{cert}{99}
\showcasevscodeiconfiletype{ceylon}{100}
\showcasevscodeiconfiletype{cf}{101}
\showcasevscodeiconfiletype{cf2}{102}
\showcasevscodeiconfiletype{cfc}{103}
\showcasevscodeiconfiletype{cfc2}{104}
\showcasevscodeiconfiletype{cfm}{105}
\showcasevscodeiconfiletype{cfm2}{106}
\showcasevscodeiconfiletype{changie}{107}
\showcasevscodeiconfiletype{cheader}{108}
\showcasevscodeiconfiletype{chef}{109}
\showcasevscodeiconfiletype{chef-cookbook}{110}
\showcasevscodeiconfiletype{circleci}{111}
\showcasevscodeiconfiletype{circom}{112}
\showcasevscodeiconfiletype{class}{113}
\showcasevscodeiconfiletype{claude}{114}
\showcasevscodeiconfiletype{clojure}{115}
\showcasevscodeiconfiletype{clojurescript}{116}
\showcasevscodeiconfiletype{cloudflare}{117}
\showcasevscodeiconfiletype{cloudfoundry}{118}
\showcasevscodeiconfiletype{cmake}{119}
\showcasevscodeiconfiletype{cobol}{120}
\showcasevscodeiconfiletype{codacy}{121}
\showcasevscodeiconfiletype{codeclimate}{122}
\showcasevscodeiconfiletype{codecov}{123}
\showcasevscodeiconfiletype{codekit}{124}
\showcasevscodeiconfiletype{codemagic}{125}
\showcasevscodeiconfiletype{codeowners}{126}
\showcasevscodeiconfiletype{codeql}{127}
\showcasevscodeiconfiletype{coderabbit}{128}
\showcasevscodeiconfiletype{coffeelint}{129}
\showcasevscodeiconfiletype{coffeescript}{130}
\showcasevscodeiconfiletype{commitizen}{131}
\showcasevscodeiconfiletype{commitlint}{132}
\showcasevscodeiconfiletype{compass}{133}
\showcasevscodeiconfiletype{composer}{134}
\showcasevscodeiconfiletype{conan}{135}
\showcasevscodeiconfiletype{conda}{136}
\showcasevscodeiconfiletype{config}{137}
\showcasevscodeiconfiletype{confluence}{138}
\showcasevscodeiconfiletype{copilot}{139}
\showcasevscodeiconfiletype{coverage}{140}
\showcasevscodeiconfiletype{coveralls}{141}
\showcasevscodeiconfiletype{cpp}{142}
\showcasevscodeiconfiletype{cpp2}{143}
\showcasevscodeiconfiletype{cpp3}{144}
\showcasevscodeiconfiletype{cppheader}{145}
\showcasevscodeiconfiletype{craco}{146}
\showcasevscodeiconfiletype{crowdin}{147}
\showcasevscodeiconfiletype{crystal}{148}
\showcasevscodeiconfiletype{csharp}{149}
\showcasevscodeiconfiletype{csharp2}{150}
\showcasevscodeiconfiletype{cspell}{151}
\showcasevscodeiconfiletype{csproj}{152}
\showcasevscodeiconfiletype{css}{153}
\showcasevscodeiconfiletype{css2}{154}
\showcasevscodeiconfiletype{css2map}{155}
\showcasevscodeiconfiletype{csscomb}{156}
\showcasevscodeiconfiletype{csslint}{157}
\showcasevscodeiconfiletype{cssmap}{158}
\showcasevscodeiconfiletype{cucumber}{159}
\showcasevscodeiconfiletype{cuda}{160}
\showcasevscodeiconfiletype{cursorrules}{161}
\showcasevscodeiconfiletype{cvs}{162}
\showcasevscodeiconfiletype{cypress}{163}
\showcasevscodeiconfiletype{cypress-spec}{164}
\showcasevscodeiconfiletype{cython}{165}
\showcasevscodeiconfiletype{c-al}{166}
\showcasevscodeiconfiletype{dal}{167}
\showcasevscodeiconfiletype{darcs}{168}
\showcasevscodeiconfiletype{dartlang}{169}
\showcasevscodeiconfiletype{dartlang-generated}{170}
\showcasevscodeiconfiletype{dartlang-ignore}{171}
\showcasevscodeiconfiletype{datadog}{172}
\showcasevscodeiconfiletype{db}{173}
\showcasevscodeiconfiletype{debian}{174}
\showcasevscodeiconfiletype{delphi}{175}
\showcasevscodeiconfiletype{deno}{176}
\showcasevscodeiconfiletype{denoify}{177}
\showcasevscodeiconfiletype{dependabot}{178}
\showcasevscodeiconfiletype{dependencies}{179}
\showcasevscodeiconfiletype{devcontainer}{180}
\showcasevscodeiconfiletype{dhall}{181}
\showcasevscodeiconfiletype{diff}{182}
\showcasevscodeiconfiletype{django}{183}
\showcasevscodeiconfiletype{dlang}{184}
\showcasevscodeiconfiletype{docker}{185}
\showcasevscodeiconfiletype{docker2}{186}
\showcasevscodeiconfiletype{dockertest}{187}
\showcasevscodeiconfiletype{dockertest2}{188}
\showcasevscodeiconfiletype{docpad}{189}
\showcasevscodeiconfiletype{docusaurus}{190}
\showcasevscodeiconfiletype{docz}{191}
\showcasevscodeiconfiletype{dojo}{192}
\showcasevscodeiconfiletype{doppler}{193}
\showcasevscodeiconfiletype{dotenv}{194}
\showcasevscodeiconfiletype{dotjs}{195}
\showcasevscodeiconfiletype{doxygen}{196}
\showcasevscodeiconfiletype{drawio}{197}
\showcasevscodeiconfiletype{drone}{198}
\showcasevscodeiconfiletype{drools}{199}
\showcasevscodeiconfiletype{dtd}{200}
\showcasevscodeiconfiletype{dune}{201}
\showcasevscodeiconfiletype{dustjs}{202}
\showcasevscodeiconfiletype{dvc}{203}
\showcasevscodeiconfiletype{dylan}{204}
\showcasevscodeiconfiletype{earthly}{205}
\showcasevscodeiconfiletype{eas-metadata}{206}
\showcasevscodeiconfiletype{edge}{207}
\showcasevscodeiconfiletype{edge2}{208}
\showcasevscodeiconfiletype{editorconfig}{209}
\showcasevscodeiconfiletype{eex}{210}
\showcasevscodeiconfiletype{ejs}{211}
\showcasevscodeiconfiletype{elastic}{212}
\showcasevscodeiconfiletype{elasticbeanstalk}{213}
\showcasevscodeiconfiletype{eleventy}{214}
\showcasevscodeiconfiletype{eleventy2}{215}
\showcasevscodeiconfiletype{elixir}{216}
\showcasevscodeiconfiletype{elm}{217}
\showcasevscodeiconfiletype{elm2}{218}
\showcasevscodeiconfiletype{emacs}{219}
\showcasevscodeiconfiletype{ember}{220}
\showcasevscodeiconfiletype{ensime}{221}
\showcasevscodeiconfiletype{eps}{222}
\showcasevscodeiconfiletype{epub}{223}
\showcasevscodeiconfiletype{erb}{224}
\showcasevscodeiconfiletype{erlang}{225}
\showcasevscodeiconfiletype{erlang2}{226}
\showcasevscodeiconfiletype{esbuild}{227}
\showcasevscodeiconfiletype{eslint}{228}
\showcasevscodeiconfiletype{eslint2}{229}
\showcasevscodeiconfiletype{esphome}{230}
\showcasevscodeiconfiletype{excalidraw}{231}
\showcasevscodeiconfiletype{excel}{232}
\showcasevscodeiconfiletype{excel2}{233}
\showcasevscodeiconfiletype{expo}{234}
\showcasevscodeiconfiletype{falcon}{235}
\showcasevscodeiconfiletype{fantasticon}{236}
\showcasevscodeiconfiletype{fastly}{237}
\showcasevscodeiconfiletype{fauna}{238}
\showcasevscodeiconfiletype{favicon}{239}
\showcasevscodeiconfiletype{fbx}{240}
\showcasevscodeiconfiletype{firebase}{241}
\showcasevscodeiconfiletype{firebasehosting}{242}
\showcasevscodeiconfiletype{firebasestorage}{243}
\showcasevscodeiconfiletype{firestore}{244}
\showcasevscodeiconfiletype{fitbit}{245}
\showcasevscodeiconfiletype{fla}{246}
\showcasevscodeiconfiletype{flareact}{247}
\showcasevscodeiconfiletype{flash}{248}
\showcasevscodeiconfiletype{floobits}{249}
\showcasevscodeiconfiletype{flow}{250}
\showcasevscodeiconfiletype{flutter}{251}
\showcasevscodeiconfiletype{flutter-package}{252}
\showcasevscodeiconfiletype{flyio}{253}
\showcasevscodeiconfiletype{font}{254}
\showcasevscodeiconfiletype{formkit}{255}
\showcasevscodeiconfiletype{fortran}{256}
\showcasevscodeiconfiletype{fossa}{257}
\showcasevscodeiconfiletype{fossil}{258}
\showcasevscodeiconfiletype{freemarker}{259}
\showcasevscodeiconfiletype{frontcommerce}{260}
\showcasevscodeiconfiletype{fsharp}{261}
\showcasevscodeiconfiletype{fsharp2}{262}
\showcasevscodeiconfiletype{fsproj}{263}
\showcasevscodeiconfiletype{fthtml}{264}
\showcasevscodeiconfiletype{funding}{265}
\showcasevscodeiconfiletype{fusebox}{266}
\showcasevscodeiconfiletype{galen}{267}
\showcasevscodeiconfiletype{galen2}{268}
\showcasevscodeiconfiletype{gamemaker}{269}
\showcasevscodeiconfiletype{gamemaker2}{270}
\showcasevscodeiconfiletype{gamemaker81}{271}
\showcasevscodeiconfiletype{gatsby}{272}
\showcasevscodeiconfiletype{gcloud}{273}
\showcasevscodeiconfiletype{gcode}{274}
\showcasevscodeiconfiletype{gdscript}{275}
\showcasevscodeiconfiletype{gemini}{276}
\showcasevscodeiconfiletype{genstat}{277}
\showcasevscodeiconfiletype{geojson}{278}
\showcasevscodeiconfiletype{git}{279}
\showcasevscodeiconfiletype{git2}{280}
\showcasevscodeiconfiletype{gitlab}{281}
\showcasevscodeiconfiletype{gitpod}{282}
\showcasevscodeiconfiletype{gleam}{283}
\showcasevscodeiconfiletype{gleamconfig}{284}
\showcasevscodeiconfiletype{glide}{285}
\showcasevscodeiconfiletype{glimmer}{286}
\showcasevscodeiconfiletype{glitter}{287}
\showcasevscodeiconfiletype{glsl}{288}
\showcasevscodeiconfiletype{gltf}{289}
\showcasevscodeiconfiletype{glyphs}{290}
\showcasevscodeiconfiletype{gnu}{291}
\showcasevscodeiconfiletype{gnuplot}{292}
\showcasevscodeiconfiletype{go}{293}
\showcasevscodeiconfiletype{goctl}{294}
\showcasevscodeiconfiletype{godot}{295}
\showcasevscodeiconfiletype{go-aqua}{296}
\showcasevscodeiconfiletype{go-black}{297}
\showcasevscodeiconfiletype{go-fuchsia}{298}
\showcasevscodeiconfiletype{go-gopher}{299}
\showcasevscodeiconfiletype{go-lightblue}{300}
\showcasevscodeiconfiletype{go-package}{301}
\showcasevscodeiconfiletype{go-white}{302}
\showcasevscodeiconfiletype{go-work}{303}
\showcasevscodeiconfiletype{go-yellow}{304}
\showcasevscodeiconfiletype{gpg}{305}
\showcasevscodeiconfiletype{gradle}{306}
\showcasevscodeiconfiletype{gradle2}{307}
\showcasevscodeiconfiletype{grain}{308}
\showcasevscodeiconfiletype{graphql}{309}
\showcasevscodeiconfiletype{graphql-config}{310}
\showcasevscodeiconfiletype{graphviz}{311}
\showcasevscodeiconfiletype{greenkeeper}{312}
\showcasevscodeiconfiletype{gridsome}{313}
\showcasevscodeiconfiletype{grit}{314}
\showcasevscodeiconfiletype{groovy}{315}
\showcasevscodeiconfiletype{groovy2}{316}
\showcasevscodeiconfiletype{grunt}{317}
\showcasevscodeiconfiletype{gulp}{318}
\showcasevscodeiconfiletype{haml}{319}
\showcasevscodeiconfiletype{handlebars}{320}
\showcasevscodeiconfiletype{handlebars2}{321}
\showcasevscodeiconfiletype{harbour}{322}
\showcasevscodeiconfiletype{hardhat}{323}
\showcasevscodeiconfiletype{hashicorp}{324}
\showcasevscodeiconfiletype{haskell}{325}
\showcasevscodeiconfiletype{haskell2}{326}
\showcasevscodeiconfiletype{haxe}{327}
\showcasevscodeiconfiletype{haxecheckstyle}{328}
\showcasevscodeiconfiletype{haxedevelop}{329}
\showcasevscodeiconfiletype{helix}{330}
\showcasevscodeiconfiletype{helm}{331}
\showcasevscodeiconfiletype{histoire}{332}
\showcasevscodeiconfiletype{hjson}{333}
\showcasevscodeiconfiletype{hlsl}{334}
\showcasevscodeiconfiletype{homeassistant}{335}
\showcasevscodeiconfiletype{horusec}{336}
\showcasevscodeiconfiletype{host}{337}
\showcasevscodeiconfiletype{html}{338}
\showcasevscodeiconfiletype{htmlhint}{339}
\showcasevscodeiconfiletype{htmlvalidate}{340}
\showcasevscodeiconfiletype{http}{341}
\showcasevscodeiconfiletype{hugo}{342}
\showcasevscodeiconfiletype{humanstxt}{343}
\showcasevscodeiconfiletype{hunspell}{344}
\showcasevscodeiconfiletype{husky}{345}
\showcasevscodeiconfiletype{hy}{346}
\showcasevscodeiconfiletype{hygen}{347}
\showcasevscodeiconfiletype{hypr}{348}
\showcasevscodeiconfiletype{icl}{349}
\showcasevscodeiconfiletype{idris}{350}
\showcasevscodeiconfiletype{idrisbin}{351}
\showcasevscodeiconfiletype{idrispkg}{352}
\showcasevscodeiconfiletype{image}{353}
\showcasevscodeiconfiletype{imba}{354}
\showcasevscodeiconfiletype{inc}{355}
\showcasevscodeiconfiletype{infopath}{356}
\showcasevscodeiconfiletype{informix}{357}
\showcasevscodeiconfiletype{ini}{358}
\showcasevscodeiconfiletype{ink}{359}
\showcasevscodeiconfiletype{innosetup}{360}
\showcasevscodeiconfiletype{io}{361}
\showcasevscodeiconfiletype{iodine}{362}
\showcasevscodeiconfiletype{ionic}{363}
\showcasevscodeiconfiletype{jake}{364}
\showcasevscodeiconfiletype{janet}{365}
\showcasevscodeiconfiletype{jar}{366}
\showcasevscodeiconfiletype{jasmine}{367}
\showcasevscodeiconfiletype{java}{368}
\showcasevscodeiconfiletype{jbuilder}{369}
\showcasevscodeiconfiletype{jekyll}{370}
\showcasevscodeiconfiletype{jenkins}{371}
\showcasevscodeiconfiletype{jest}{372}
\showcasevscodeiconfiletype{jest-snapshot}{373}
\showcasevscodeiconfiletype{jinja}{374}
\showcasevscodeiconfiletype{jpm}{375}
\showcasevscodeiconfiletype{js}{376}
\showcasevscodeiconfiletype{jsbeautify}{377}
\showcasevscodeiconfiletype{jsconfig}{378}
\showcasevscodeiconfiletype{jscpd}{379}
\showcasevscodeiconfiletype{jshint}{380}
\showcasevscodeiconfiletype{jsmap}{381}
\showcasevscodeiconfiletype{json}{382}
\showcasevscodeiconfiletype{json2}{383}
\showcasevscodeiconfiletype{json5}{384}
\showcasevscodeiconfiletype{jsonld}{385}
\showcasevscodeiconfiletype{jsonnet}{386}
\showcasevscodeiconfiletype{json-official}{387}
\showcasevscodeiconfiletype{json-schema}{388}
\showcasevscodeiconfiletype{jsp}{389}
\showcasevscodeiconfiletype{jsr}{390}
\showcasevscodeiconfiletype{jss}{391}
\showcasevscodeiconfiletype{js-official}{392}
\showcasevscodeiconfiletype{juice}{393}
\showcasevscodeiconfiletype{julia}{394}
\showcasevscodeiconfiletype{julia2}{395}
\showcasevscodeiconfiletype{jupyter}{396}
\showcasevscodeiconfiletype{just}{397}
\showcasevscodeiconfiletype{k}{398}
\showcasevscodeiconfiletype{karma}{399}
\showcasevscodeiconfiletype{key}{400}
\showcasevscodeiconfiletype{kitchenci}{401}
\showcasevscodeiconfiletype{kite}{402}
\showcasevscodeiconfiletype{kivy}{403}
\showcasevscodeiconfiletype{knip}{404}
\showcasevscodeiconfiletype{kos}{405}
\showcasevscodeiconfiletype{kotlin}{406}
\showcasevscodeiconfiletype{kusto}{407}
\showcasevscodeiconfiletype{language-configuration}{408}
\showcasevscodeiconfiletype{lark}{409}
\showcasevscodeiconfiletype{latino}{410}
\showcasevscodeiconfiletype{layout}{411}
\showcasevscodeiconfiletype{lean}{412}
\showcasevscodeiconfiletype{leanconfig}{413}
\showcasevscodeiconfiletype{lefthook}{414}
\showcasevscodeiconfiletype{lerna}{415}
\showcasevscodeiconfiletype{less}{416}
\showcasevscodeiconfiletype{lex}{417}
\showcasevscodeiconfiletype{liara}{418}
\showcasevscodeiconfiletype{libreoffice-base}{419}
\showcasevscodeiconfiletype{libreoffice-calc}{420}
\showcasevscodeiconfiletype{libreoffice-draw}{421}
\showcasevscodeiconfiletype{libreoffice-impress}{422}
\showcasevscodeiconfiletype{libreoffice-math}{423}
\showcasevscodeiconfiletype{libreoffice-writer}{424}
\showcasevscodeiconfiletype{license}{425}
\showcasevscodeiconfiletype{licensebat}{426}
\showcasevscodeiconfiletype{lighthouse}{427}
\showcasevscodeiconfiletype{light-actionscript2}{428}
\showcasevscodeiconfiletype{light-ada}{429}
\showcasevscodeiconfiletype{light-agda}{430}
\showcasevscodeiconfiletype{light-agents}{431}
\showcasevscodeiconfiletype{light-apl}{432}
\showcasevscodeiconfiletype{light-astro}{433}
\showcasevscodeiconfiletype{light-astroconfig}{434}
\showcasevscodeiconfiletype{light-babel}{435}
\showcasevscodeiconfiletype{light-babel2}{436}
\showcasevscodeiconfiletype{light-cabal}{437}
\showcasevscodeiconfiletype{light-circleci}{438}
\showcasevscodeiconfiletype{light-cloudfoundry}{439}
\showcasevscodeiconfiletype{light-codacy}{440}
\showcasevscodeiconfiletype{light-codeclimate}{441}
\showcasevscodeiconfiletype{light-codeowners}{442}
\showcasevscodeiconfiletype{light-config}{443}
\showcasevscodeiconfiletype{light-copilot}{444}
\showcasevscodeiconfiletype{light-crystal}{445}
\showcasevscodeiconfiletype{light-cypress}{446}
\showcasevscodeiconfiletype{light-cypress-spec}{447}
\showcasevscodeiconfiletype{light-db}{448}
\showcasevscodeiconfiletype{light-dhall}{449}
\showcasevscodeiconfiletype{light-docpad}{450}
\showcasevscodeiconfiletype{light-drone}{451}
\showcasevscodeiconfiletype{light-eas-metadata}{452}
\showcasevscodeiconfiletype{light-eleventy}{453}
\showcasevscodeiconfiletype{light-eleventy2}{454}
\showcasevscodeiconfiletype{light-esphome}{455}
\showcasevscodeiconfiletype{light-expo}{456}
\showcasevscodeiconfiletype{light-firebasehosting}{457}
\showcasevscodeiconfiletype{light-fla}{458}
\showcasevscodeiconfiletype{light-font}{459}
\showcasevscodeiconfiletype{light-gamemaker2}{460}
\showcasevscodeiconfiletype{light-gradle}{461}
\showcasevscodeiconfiletype{light-hashicorp}{462}
\showcasevscodeiconfiletype{light-hjson}{463}
\showcasevscodeiconfiletype{light-ini}{464}
\showcasevscodeiconfiletype{light-io}{465}
\showcasevscodeiconfiletype{light-js}{466}
\showcasevscodeiconfiletype{light-jsconfig}{467}
\showcasevscodeiconfiletype{light-jsmap}{468}
\showcasevscodeiconfiletype{light-json}{469}
\showcasevscodeiconfiletype{light-json5}{470}
\showcasevscodeiconfiletype{light-jsonld}{471}
\showcasevscodeiconfiletype{light-json-schema}{472}
\showcasevscodeiconfiletype{light-kite}{473}
\showcasevscodeiconfiletype{light-lark}{474}
\showcasevscodeiconfiletype{light-lean}{475}
\showcasevscodeiconfiletype{light-leanconfig}{476}
\showcasevscodeiconfiletype{light-lerna}{477}
\showcasevscodeiconfiletype{light-mailing}{478}
\showcasevscodeiconfiletype{light-markuplint}{479}
\showcasevscodeiconfiletype{light-mcp}{480}
\showcasevscodeiconfiletype{light-mdx-components}{481}
\showcasevscodeiconfiletype{light-mdx}{482}
\showcasevscodeiconfiletype{light-mlang}{483}
\showcasevscodeiconfiletype{light-mustache}{484}
\showcasevscodeiconfiletype{light-mypy}{485}
\showcasevscodeiconfiletype{light-neo4j}{486}
\showcasevscodeiconfiletype{light-netlify}{487}
\showcasevscodeiconfiletype{light-next}{488}
\showcasevscodeiconfiletype{light-nim}{489}
\showcasevscodeiconfiletype{light-nx}{490}
\showcasevscodeiconfiletype{light-objidconfig}{491}
\showcasevscodeiconfiletype{light-openHAB}{492}
\showcasevscodeiconfiletype{light-packship}{493}
\showcasevscodeiconfiletype{light-pcl}{494}
\showcasevscodeiconfiletype{light-pnpm}{495}
\showcasevscodeiconfiletype{light-prettier}{496}
\showcasevscodeiconfiletype{light-prisma}{497}
\showcasevscodeiconfiletype{light-purescript}{498}
\showcasevscodeiconfiletype{light-quasar}{499}
\showcasevscodeiconfiletype{light-raku}{500}
\showcasevscodeiconfiletype{light-razzle}{501}
\showcasevscodeiconfiletype{light-reactrouter}{502}
\showcasevscodeiconfiletype{light-rehype}{503}
\showcasevscodeiconfiletype{light-remark}{504}
\showcasevscodeiconfiletype{light-replit}{505}
\showcasevscodeiconfiletype{light-retext}{506}
\showcasevscodeiconfiletype{light-rome}{507}
\showcasevscodeiconfiletype{light-rubocop}{508}
\showcasevscodeiconfiletype{light-rust}{509}
\showcasevscodeiconfiletype{light-rust-toolchain}{510}
\showcasevscodeiconfiletype{light-safetensors}{511}
\showcasevscodeiconfiletype{light-shadcn}{512}
\showcasevscodeiconfiletype{light-shaderlab}{513}
\showcasevscodeiconfiletype{light-solidity}{514}
\showcasevscodeiconfiletype{light-spin}{515}
\showcasevscodeiconfiletype{light-stylelint}{516}
\showcasevscodeiconfiletype{light-stylus}{517}
\showcasevscodeiconfiletype{light-symfony}{518}
\showcasevscodeiconfiletype{light-systemd}{519}
\showcasevscodeiconfiletype{light-systemverilog}{520}
\showcasevscodeiconfiletype{light-testcafe}{521}
\showcasevscodeiconfiletype{light-testjs}{522}
\showcasevscodeiconfiletype{light-tex}{523}
\showcasevscodeiconfiletype{light-tm}{524}
\showcasevscodeiconfiletype{light-tmux}{525}
\showcasevscodeiconfiletype{light-todo}{526}
\showcasevscodeiconfiletype{light-toit}{527}
\showcasevscodeiconfiletype{light-toml}{528}
\showcasevscodeiconfiletype{light-tree}{529}
\showcasevscodeiconfiletype{light-turbo}{530}
\showcasevscodeiconfiletype{light-unibeautify}{531}
\showcasevscodeiconfiletype{light-vash}{532}
\showcasevscodeiconfiletype{light-vercel}{533}
\showcasevscodeiconfiletype{light-vsix}{534}
\showcasevscodeiconfiletype{light-vsixmanifest}{535}
\showcasevscodeiconfiletype{light-xfl}{536}
\showcasevscodeiconfiletype{light-xorg}{537}
\showcasevscodeiconfiletype{light-yaml}{538}
\showcasevscodeiconfiletype{light-yaml-official}{539}
\showcasevscodeiconfiletype{lilypond}{540}
\showcasevscodeiconfiletype{lime}{541}
\showcasevscodeiconfiletype{lintstagedrc}{542}
\showcasevscodeiconfiletype{liquid}{543}
\showcasevscodeiconfiletype{lisp}{544}
\showcasevscodeiconfiletype{livescript}{545}
\showcasevscodeiconfiletype{lnk}{546}
\showcasevscodeiconfiletype{locale}{547}
\showcasevscodeiconfiletype{log}{548}
\showcasevscodeiconfiletype{lolcode}{549}
\showcasevscodeiconfiletype{lsl}{550}
\showcasevscodeiconfiletype{lua}{551}
\showcasevscodeiconfiletype{luau}{552}
\showcasevscodeiconfiletype{lync}{553}
\showcasevscodeiconfiletype{mailing}{554}
\showcasevscodeiconfiletype{manifest}{555}
\showcasevscodeiconfiletype{manifest-bak}{556}
\showcasevscodeiconfiletype{manifest-skip}{557}
\showcasevscodeiconfiletype{map}{558}
\showcasevscodeiconfiletype{mariadb}{559}
\showcasevscodeiconfiletype{markdown}{560}
\showcasevscodeiconfiletype{markdownlint}{561}
\showcasevscodeiconfiletype{markdownlint-ignore}{562}
\showcasevscodeiconfiletype{marko}{563}
\showcasevscodeiconfiletype{markojs}{564}
\showcasevscodeiconfiletype{markuplint}{565}
\showcasevscodeiconfiletype{master-co}{566}
\showcasevscodeiconfiletype{matlab}{567}
\showcasevscodeiconfiletype{maven}{568}
\showcasevscodeiconfiletype{maxscript}{569}
\showcasevscodeiconfiletype{maya}{570}
\showcasevscodeiconfiletype{mcp}{571}
\showcasevscodeiconfiletype{mdx-components}{572}
\showcasevscodeiconfiletype{mdx}{573}
\showcasevscodeiconfiletype{mdxlint}{574}
\showcasevscodeiconfiletype{mediawiki}{575}
\showcasevscodeiconfiletype{mercurial}{576}
\showcasevscodeiconfiletype{mermaid}{577}
\showcasevscodeiconfiletype{meson}{578}
\showcasevscodeiconfiletype{metal}{579}
\showcasevscodeiconfiletype{meteor}{580}
\showcasevscodeiconfiletype{minecraft}{581}
\showcasevscodeiconfiletype{mivascript}{582}
\showcasevscodeiconfiletype{mjml}{583}
\showcasevscodeiconfiletype{mlang}{584}
\showcasevscodeiconfiletype{mocha}{585}
\showcasevscodeiconfiletype{modernizr}{586}
\showcasevscodeiconfiletype{mojo}{587}
\showcasevscodeiconfiletype{mojolicious}{588}
\showcasevscodeiconfiletype{moleculer}{589}
\showcasevscodeiconfiletype{mondoo}{590}
\showcasevscodeiconfiletype{mongo}{591}
\showcasevscodeiconfiletype{monotone}{592}
\showcasevscodeiconfiletype{motif}{593}
\showcasevscodeiconfiletype{mson}{594}
\showcasevscodeiconfiletype{mustache}{595}
\showcasevscodeiconfiletype{mvt}{596}
\showcasevscodeiconfiletype{mvtcss}{597}
\showcasevscodeiconfiletype{mvtjs}{598}
\showcasevscodeiconfiletype{mypy}{599}
\showcasevscodeiconfiletype{mysql}{600}
\showcasevscodeiconfiletype{nanostaged}{601}
\showcasevscodeiconfiletype{ndst}{602}
\showcasevscodeiconfiletype{nearly}{603}
\showcasevscodeiconfiletype{neo4j}{604}
\showcasevscodeiconfiletype{nestjs}{605}
\showcasevscodeiconfiletype{nest-adapter-js}{606}
\showcasevscodeiconfiletype{nest-adapter-ts}{607}
\showcasevscodeiconfiletype{nest-controller-js}{608}
\showcasevscodeiconfiletype{nest-controller-ts}{609}
\showcasevscodeiconfiletype{nest-decorator-js}{610}
\showcasevscodeiconfiletype{nest-decorator-ts}{611}
\showcasevscodeiconfiletype{nest-filter-js}{612}
\showcasevscodeiconfiletype{nest-filter-ts}{613}
\showcasevscodeiconfiletype{nest-gateway-js}{614}
\showcasevscodeiconfiletype{nest-gateway-ts}{615}
\showcasevscodeiconfiletype{nest-guard-js}{616}
\showcasevscodeiconfiletype{nest-guard-ts}{617}
\showcasevscodeiconfiletype{nest-interceptor-js}{618}
\showcasevscodeiconfiletype{nest-interceptor-ts}{619}
\showcasevscodeiconfiletype{nest-middleware-js}{620}
\showcasevscodeiconfiletype{nest-middleware-ts}{621}
\showcasevscodeiconfiletype{nest-module-js}{622}
\showcasevscodeiconfiletype{nest-module-ts}{623}
\showcasevscodeiconfiletype{nest-pipe-js}{624}
\showcasevscodeiconfiletype{nest-pipe-ts}{625}
\showcasevscodeiconfiletype{nest-service-js}{626}
\showcasevscodeiconfiletype{nest-service-ts}{627}
\showcasevscodeiconfiletype{netlify}{628}
\showcasevscodeiconfiletype{next}{629}
\showcasevscodeiconfiletype{nextflow}{630}
\showcasevscodeiconfiletype{nginx}{631}
\showcasevscodeiconfiletype{ng-component-css}{632}
\showcasevscodeiconfiletype{ng-component-dart}{633}
\showcasevscodeiconfiletype{ng-component-html}{634}
\showcasevscodeiconfiletype{ng-component-js}{635}
\showcasevscodeiconfiletype{ng-component-js2}{636}
\showcasevscodeiconfiletype{ng-component-less}{637}
\showcasevscodeiconfiletype{ng-component-sass}{638}
\showcasevscodeiconfiletype{ng-component-scss}{639}
\showcasevscodeiconfiletype{ng-component-ts}{640}
\showcasevscodeiconfiletype{ng-component-ts2}{641}
\showcasevscodeiconfiletype{ng-controller-js}{642}
\showcasevscodeiconfiletype{ng-controller-ts}{643}
\showcasevscodeiconfiletype{ng-directive-dart}{644}
\showcasevscodeiconfiletype{ng-directive-js}{645}
\showcasevscodeiconfiletype{ng-directive-js2}{646}
\showcasevscodeiconfiletype{ng-directive-ts}{647}
\showcasevscodeiconfiletype{ng-directive-ts2}{648}
\showcasevscodeiconfiletype{ng-guard-dart}{649}
\showcasevscodeiconfiletype{ng-guard-js}{650}
\showcasevscodeiconfiletype{ng-guard-ts}{651}
\showcasevscodeiconfiletype{ng-interceptor-dart}{652}
\showcasevscodeiconfiletype{ng-interceptor-js}{653}
\showcasevscodeiconfiletype{ng-interceptor-ts}{654}
\showcasevscodeiconfiletype{ng-module-dart}{655}
\showcasevscodeiconfiletype{ng-module-js}{656}
\showcasevscodeiconfiletype{ng-module-js2}{657}
\showcasevscodeiconfiletype{ng-module-ts}{658}
\showcasevscodeiconfiletype{ng-module-ts2}{659}
\showcasevscodeiconfiletype{ng-pipe-dart}{660}
\showcasevscodeiconfiletype{ng-pipe-js}{661}
\showcasevscodeiconfiletype{ng-pipe-js2}{662}
\showcasevscodeiconfiletype{ng-pipe-ts}{663}
\showcasevscodeiconfiletype{ng-pipe-ts2}{664}
\showcasevscodeiconfiletype{ng-routing-dart}{665}
\showcasevscodeiconfiletype{ng-routing-js}{666}
\showcasevscodeiconfiletype{ng-routing-js2}{667}
\showcasevscodeiconfiletype{ng-routing-ts}{668}
\showcasevscodeiconfiletype{ng-routing-ts2}{669}
\showcasevscodeiconfiletype{ng-service-dart}{670}
\showcasevscodeiconfiletype{ng-service-js}{671}
\showcasevscodeiconfiletype{ng-service-js2}{672}
\showcasevscodeiconfiletype{ng-service-ts}{673}
\showcasevscodeiconfiletype{ng-service-ts2}{674}
\showcasevscodeiconfiletype{ng-smart-component-dart}{675}
\showcasevscodeiconfiletype{ng-smart-component-js}{676}
\showcasevscodeiconfiletype{ng-smart-component-js2}{677}
\showcasevscodeiconfiletype{ng-smart-component-ts}{678}
\showcasevscodeiconfiletype{ng-smart-component-ts2}{679}
\showcasevscodeiconfiletype{ng-tailwind}{680}
\showcasevscodeiconfiletype{nim}{681}
\showcasevscodeiconfiletype{nimble}{682}
\showcasevscodeiconfiletype{ninja}{683}
\showcasevscodeiconfiletype{nix}{684}
\showcasevscodeiconfiletype{njsproj}{685}
\showcasevscodeiconfiletype{noc}{686}
\showcasevscodeiconfiletype{node}{687}
\showcasevscodeiconfiletype{node2}{688}
\showcasevscodeiconfiletype{nodemon}{689}
\showcasevscodeiconfiletype{npm}{690}
\showcasevscodeiconfiletype{npmpackagejsonlint}{691}
\showcasevscodeiconfiletype{nsi}{692}
\showcasevscodeiconfiletype{nsri-integrity}{693}
\showcasevscodeiconfiletype{nsri}{694}
\showcasevscodeiconfiletype{nuget}{695}
\showcasevscodeiconfiletype{numpy}{696}
\showcasevscodeiconfiletype{nunjucks}{697}
\showcasevscodeiconfiletype{nuxt}{698}
\showcasevscodeiconfiletype{nx}{699}
\showcasevscodeiconfiletype{nyc}{700}
\showcasevscodeiconfiletype{objectivec}{701}
\showcasevscodeiconfiletype{objectivecpp}{702}
\showcasevscodeiconfiletype{objidconfig}{703}
\showcasevscodeiconfiletype{ocaml}{704}
\showcasevscodeiconfiletype{ocaml-intf}{705}
\showcasevscodeiconfiletype{ogone}{706}
\showcasevscodeiconfiletype{onenote}{707}
\showcasevscodeiconfiletype{opam}{708}
\showcasevscodeiconfiletype{opencl}{709}
\showcasevscodeiconfiletype{openHAB}{710}
\showcasevscodeiconfiletype{openscad}{711}
\showcasevscodeiconfiletype{opentofu}{712}
\showcasevscodeiconfiletype{org}{713}
\showcasevscodeiconfiletype{outlook}{714}
\showcasevscodeiconfiletype{ovpn}{715}
\showcasevscodeiconfiletype{oxc}{716}
\showcasevscodeiconfiletype{package}{717}
\showcasevscodeiconfiletype{packship}{718}
\showcasevscodeiconfiletype{paket}{719}
\showcasevscodeiconfiletype{pandacss}{720}
\showcasevscodeiconfiletype{patch}{721}
\showcasevscodeiconfiletype{pcl}{722}
\showcasevscodeiconfiletype{pddl}{723}
\showcasevscodeiconfiletype{pddl-happenings}{724}
\showcasevscodeiconfiletype{pddl-plan}{725}
\showcasevscodeiconfiletype{pdf}{726}
\showcasevscodeiconfiletype{pdf2}{727}
\showcasevscodeiconfiletype{pdm}{728}
\showcasevscodeiconfiletype{peeky}{729}
\showcasevscodeiconfiletype{perl}{730}
\showcasevscodeiconfiletype{perl2}{731}
\showcasevscodeiconfiletype{perl6}{732}
\showcasevscodeiconfiletype{pgsql}{733}
\showcasevscodeiconfiletype{photoshop}{734}
\showcasevscodeiconfiletype{photoshop2}{735}
\showcasevscodeiconfiletype{php}{736}
\showcasevscodeiconfiletype{php2}{737}
\showcasevscodeiconfiletype{php3}{738}
\showcasevscodeiconfiletype{phpcsfixer}{739}
\showcasevscodeiconfiletype{phpstan}{740}
\showcasevscodeiconfiletype{phpunit}{741}
\showcasevscodeiconfiletype{phraseapp}{742}
\showcasevscodeiconfiletype{pine}{743}
\showcasevscodeiconfiletype{pip}{744}
\showcasevscodeiconfiletype{pipeline}{745}
\showcasevscodeiconfiletype{plantuml}{746}
\showcasevscodeiconfiletype{platformio}{747}
\showcasevscodeiconfiletype{playwright}{748}
\showcasevscodeiconfiletype{plsql}{749}
\showcasevscodeiconfiletype{plsql-package}{750}
\showcasevscodeiconfiletype{plsql-package-body}{751}
\showcasevscodeiconfiletype{plsql-package-header}{752}
\showcasevscodeiconfiletype{plsql-package-spec}{753}
\showcasevscodeiconfiletype{pm2}{754}
\showcasevscodeiconfiletype{pnpm}{755}
\showcasevscodeiconfiletype{poedit}{756}
\showcasevscodeiconfiletype{poetry}{757}
\showcasevscodeiconfiletype{polymer}{758}
\showcasevscodeiconfiletype{pony}{759}
\showcasevscodeiconfiletype{postcss}{760}
\showcasevscodeiconfiletype{postcssconfig}{761}
\showcasevscodeiconfiletype{postman}{762}
\showcasevscodeiconfiletype{powerpoint}{763}
\showcasevscodeiconfiletype{powerpoint2}{764}
\showcasevscodeiconfiletype{powershell}{765}
\showcasevscodeiconfiletype{powershell2}{766}
\showcasevscodeiconfiletype{powershell-format}{767}
\showcasevscodeiconfiletype{powershell-psd}{768}
\showcasevscodeiconfiletype{powershell-psd2}{769}
\showcasevscodeiconfiletype{powershell-psm}{770}
\showcasevscodeiconfiletype{powershell-psm2}{771}
\showcasevscodeiconfiletype{powershell-types}{772}
\showcasevscodeiconfiletype{preact}{773}
\showcasevscodeiconfiletype{precommit}{774}
\showcasevscodeiconfiletype{prettier}{775}
\showcasevscodeiconfiletype{prisma}{776}
\showcasevscodeiconfiletype{processinglang}{777}
\showcasevscodeiconfiletype{procfile}{778}
\showcasevscodeiconfiletype{progress}{779}
\showcasevscodeiconfiletype{prolog}{780}
\showcasevscodeiconfiletype{prometheus}{781}
\showcasevscodeiconfiletype{protobuf}{782}
\showcasevscodeiconfiletype{protractor}{783}
\showcasevscodeiconfiletype{publiccode}{784}
\showcasevscodeiconfiletype{publisher}{785}
\showcasevscodeiconfiletype{pug}{786}
\showcasevscodeiconfiletype{pulumi}{787}
\showcasevscodeiconfiletype{puppet}{788}
\showcasevscodeiconfiletype{purescript}{789}
\showcasevscodeiconfiletype{purgecss}{790}
\showcasevscodeiconfiletype{pyenv}{791}
\showcasevscodeiconfiletype{pyret}{792}
\showcasevscodeiconfiletype{pyscript}{793}
\showcasevscodeiconfiletype{pytest}{794}
\showcasevscodeiconfiletype{python}{795}
\showcasevscodeiconfiletype{pythowo}{796}
\showcasevscodeiconfiletype{pytyped}{797}
\showcasevscodeiconfiletype{pyup}{798}
\showcasevscodeiconfiletype{q}{799}
\showcasevscodeiconfiletype{qbs}{800}
\showcasevscodeiconfiletype{qlikview}{801}
\showcasevscodeiconfiletype{qml}{802}
\showcasevscodeiconfiletype{qmldir}{803}
\showcasevscodeiconfiletype{qsharp}{804}
\showcasevscodeiconfiletype{quasar}{805}
\showcasevscodeiconfiletype{r}{806}
\showcasevscodeiconfiletype{racket}{807}
\showcasevscodeiconfiletype{rails}{808}
\showcasevscodeiconfiletype{rake}{809}
\showcasevscodeiconfiletype{raku}{810}
\showcasevscodeiconfiletype{raml}{811}
\showcasevscodeiconfiletype{razor}{812}
\showcasevscodeiconfiletype{razzle}{813}
\showcasevscodeiconfiletype{ra-syntax-tree}{814}
\showcasevscodeiconfiletype{reactjs}{815}
\showcasevscodeiconfiletype{reactrouter}{816}
\showcasevscodeiconfiletype{reacttemplate}{817}
\showcasevscodeiconfiletype{reactts}{818}
\showcasevscodeiconfiletype{reason}{819}
\showcasevscodeiconfiletype{red}{820}
\showcasevscodeiconfiletype{registry}{821}
\showcasevscodeiconfiletype{rego}{822}
\showcasevscodeiconfiletype{rehype}{823}
\showcasevscodeiconfiletype{remark}{824}
\showcasevscodeiconfiletype{renovate}{825}
\showcasevscodeiconfiletype{replit}{826}
\showcasevscodeiconfiletype{rescript}{827}
\showcasevscodeiconfiletype{rest}{828}
\showcasevscodeiconfiletype{retext}{829}
\showcasevscodeiconfiletype{rexx}{830}
\showcasevscodeiconfiletype{riot}{831}
\showcasevscodeiconfiletype{ripple}{832}
\showcasevscodeiconfiletype{rmd}{833}
\showcasevscodeiconfiletype{rnc}{834}
\showcasevscodeiconfiletype{robotframework}{835}
\showcasevscodeiconfiletype{robots}{836}
\showcasevscodeiconfiletype{rolldown}{837}
\showcasevscodeiconfiletype{rollup}{838}
\showcasevscodeiconfiletype{rome}{839}
\showcasevscodeiconfiletype{ron}{840}
\showcasevscodeiconfiletype{rproj}{841}
\showcasevscodeiconfiletype{rspec}{842}
\showcasevscodeiconfiletype{rss}{843}
\showcasevscodeiconfiletype{rubocop}{844}
\showcasevscodeiconfiletype{ruby}{845}
\showcasevscodeiconfiletype{rust}{846}
\showcasevscodeiconfiletype{rust-toolchain}{847}
\showcasevscodeiconfiletype{s-lang}{848}
\showcasevscodeiconfiletype{safetensors}{849}
\showcasevscodeiconfiletype{sails}{850}
\showcasevscodeiconfiletype{saltstack}{851}
\showcasevscodeiconfiletype{san}{852}
\showcasevscodeiconfiletype{sapphire-framework-cli}{853}
\showcasevscodeiconfiletype{sas}{854}
\showcasevscodeiconfiletype{sass}{855}
\showcasevscodeiconfiletype{sbt}{856}
\showcasevscodeiconfiletype{scala}{857}
\showcasevscodeiconfiletype{scilab}{858}
\showcasevscodeiconfiletype{script}{859}
\showcasevscodeiconfiletype{scss}{860}
\showcasevscodeiconfiletype{scss2}{861}
\showcasevscodeiconfiletype{sdlang}{862}
\showcasevscodeiconfiletype{search-result}{863}
\showcasevscodeiconfiletype{seedkit}{864}
\showcasevscodeiconfiletype{sentry}{865}
\showcasevscodeiconfiletype{sequelize}{866}
\showcasevscodeiconfiletype{serverless}{867}
\showcasevscodeiconfiletype{shadcn}{868}
\showcasevscodeiconfiletype{shaderlab}{869}
\showcasevscodeiconfiletype{shell}{870}
\showcasevscodeiconfiletype{shellcheck}{871}
\showcasevscodeiconfiletype{shuttle}{872}
\showcasevscodeiconfiletype{silverstripe}{873}
\showcasevscodeiconfiletype{sino}{874}
\showcasevscodeiconfiletype{siyuan}{875}
\showcasevscodeiconfiletype{sketch}{876}
\showcasevscodeiconfiletype{skipper}{877}
\showcasevscodeiconfiletype{slang}{878}
\showcasevscodeiconfiletype{slashup}{879}
\showcasevscodeiconfiletype{slice}{880}
\showcasevscodeiconfiletype{slim}{881}
\showcasevscodeiconfiletype{slint}{882}
\showcasevscodeiconfiletype{sln}{883}
\showcasevscodeiconfiletype{sln2}{884}
\showcasevscodeiconfiletype{smarty}{885}
\showcasevscodeiconfiletype{smithery}{886}
\showcasevscodeiconfiletype{snakemake}{887}
\showcasevscodeiconfiletype{snapcraft}{888}
\showcasevscodeiconfiletype{snaplet}{889}
\showcasevscodeiconfiletype{snort}{890}
\showcasevscodeiconfiletype{snyk}{891}
\showcasevscodeiconfiletype{solidarity}{892}
\showcasevscodeiconfiletype{solidity}{893}
\showcasevscodeiconfiletype{source}{894}
\showcasevscodeiconfiletype{spacengine}{895}
\showcasevscodeiconfiletype{sparql}{896}
\showcasevscodeiconfiletype{spin}{897}
\showcasevscodeiconfiletype{sqf}{898}
\showcasevscodeiconfiletype{sql}{899}
\showcasevscodeiconfiletype{sqlite}{900}
\showcasevscodeiconfiletype{squirrel}{901}
\showcasevscodeiconfiletype{sss}{902}
\showcasevscodeiconfiletype{sst}{903}
\showcasevscodeiconfiletype{stackblitz}{904}
\showcasevscodeiconfiletype{stan}{905}
\showcasevscodeiconfiletype{stata}{906}
\showcasevscodeiconfiletype{stencil}{907}
\showcasevscodeiconfiletype{storyboard}{908}
\showcasevscodeiconfiletype{storybook}{909}
\showcasevscodeiconfiletype{stryker}{910}
\showcasevscodeiconfiletype{stylable}{911}
\showcasevscodeiconfiletype{style}{912}
\showcasevscodeiconfiletype{styled}{913}
\showcasevscodeiconfiletype{stylelint}{914}
\showcasevscodeiconfiletype{stylish-haskell}{915}
\showcasevscodeiconfiletype{stylus}{916}
\showcasevscodeiconfiletype{sublime}{917}
\showcasevscodeiconfiletype{subversion}{918}
\showcasevscodeiconfiletype{svelte}{919}
\showcasevscodeiconfiletype{svelteconfig}{920}
\showcasevscodeiconfiletype{svg}{921}
\showcasevscodeiconfiletype{svgo}{922}
\showcasevscodeiconfiletype{swagger}{923}
\showcasevscodeiconfiletype{swc}{924}
\showcasevscodeiconfiletype{swift}{925}
\showcasevscodeiconfiletype{swig}{926}
\showcasevscodeiconfiletype{symfony}{927}
\showcasevscodeiconfiletype{syncpack}{928}
\showcasevscodeiconfiletype{systemd}{929}
\showcasevscodeiconfiletype{systemverilog}{930}
\showcasevscodeiconfiletype{t4tt}{931}
\showcasevscodeiconfiletype{tailwind}{932}
\showcasevscodeiconfiletype{tamagui}{933}
\showcasevscodeiconfiletype{taplo}{934}
\showcasevscodeiconfiletype{taskfile}{935}
\showcasevscodeiconfiletype{tauri}{936}
\showcasevscodeiconfiletype{tcl}{937}
\showcasevscodeiconfiletype{teal}{938}
\showcasevscodeiconfiletype{templ}{939}
\showcasevscodeiconfiletype{tera}{940}
\showcasevscodeiconfiletype{terraform}{941}
\showcasevscodeiconfiletype{test}{942}
\showcasevscodeiconfiletype{testcafe}{943}
\showcasevscodeiconfiletype{testjs}{944}
\showcasevscodeiconfiletype{testplane}{945}
\showcasevscodeiconfiletype{testts}{946}
\showcasevscodeiconfiletype{tex}{947}
\showcasevscodeiconfiletype{text}{948}
\showcasevscodeiconfiletype{textile}{949}
\showcasevscodeiconfiletype{tfs}{950}
\showcasevscodeiconfiletype{tiltfile}{951}
\showcasevscodeiconfiletype{tm}{952}
\showcasevscodeiconfiletype{tmux}{953}
\showcasevscodeiconfiletype{todo}{954}
\showcasevscodeiconfiletype{toit}{955}
\showcasevscodeiconfiletype{toml}{956}
\showcasevscodeiconfiletype{tox}{957}
\showcasevscodeiconfiletype{travis}{958}
\showcasevscodeiconfiletype{tree}{959}
\showcasevscodeiconfiletype{tres}{960}
\showcasevscodeiconfiletype{truffle}{961}
\showcasevscodeiconfiletype{trunk}{962}
\showcasevscodeiconfiletype{tsbuildinfo}{963}
\showcasevscodeiconfiletype{tscn}{964}
\showcasevscodeiconfiletype{tsconfig}{965}
\showcasevscodeiconfiletype{tsconfig-official}{966}
\showcasevscodeiconfiletype{tsdoc}{967}
\showcasevscodeiconfiletype{tsdown}{968}
\showcasevscodeiconfiletype{tslint}{969}
\showcasevscodeiconfiletype{tt}{970}
\showcasevscodeiconfiletype{ttcn}{971}
\showcasevscodeiconfiletype{tuc}{972}
\showcasevscodeiconfiletype{turbo}{973}
\showcasevscodeiconfiletype{twig}{974}
\showcasevscodeiconfiletype{typedoc}{975}
\showcasevscodeiconfiletype{typescript}{976}
\showcasevscodeiconfiletype{typescriptdef}{977}
\showcasevscodeiconfiletype{typescriptdef-official}{978}
\showcasevscodeiconfiletype{typescript-official}{979}
\showcasevscodeiconfiletype{typo3}{980}
\showcasevscodeiconfiletype{uiua}{981}
\showcasevscodeiconfiletype{unibeautify}{982}
\showcasevscodeiconfiletype{unison}{983}
\showcasevscodeiconfiletype{unlicense}{984}
\showcasevscodeiconfiletype{unocss}{985}
\showcasevscodeiconfiletype{uv}{986}
\showcasevscodeiconfiletype{vagrant}{987}
\showcasevscodeiconfiletype{vala}{988}
\showcasevscodeiconfiletype{vanilla-extract}{989}
\showcasevscodeiconfiletype{vapi}{990}
\showcasevscodeiconfiletype{vapor}{991}
\showcasevscodeiconfiletype{vash}{992}
\showcasevscodeiconfiletype{vb}{993}
\showcasevscodeiconfiletype{vba}{994}
\showcasevscodeiconfiletype{vbhtml}{995}
\showcasevscodeiconfiletype{vbproj}{996}
\showcasevscodeiconfiletype{vcxproj}{997}
\showcasevscodeiconfiletype{velocity}{998}
\showcasevscodeiconfiletype{vento}{999}
\showcasevscodeiconfiletype{vercel}{1000}
\showcasevscodeiconfiletype{verilog}{1001}
\showcasevscodeiconfiletype{vhdl}{1002}
\showcasevscodeiconfiletype{video}{1003}
\showcasevscodeiconfiletype{view}{1004}
\showcasevscodeiconfiletype{vim}{1005}
\showcasevscodeiconfiletype{vite}{1006}
\showcasevscodeiconfiletype{vitest}{1007}
\showcasevscodeiconfiletype{vlang}{1008}
\showcasevscodeiconfiletype{volt}{1009}
\showcasevscodeiconfiletype{vscode-insiders}{1010}
\showcasevscodeiconfiletype{vscode}{1011}
\showcasevscodeiconfiletype{vscode2}{1012}
\showcasevscodeiconfiletype{vscode3}{1013}
\showcasevscodeiconfiletype{vscode-test}{1014}
\showcasevscodeiconfiletype{vsix}{1015}
\showcasevscodeiconfiletype{vsixmanifest}{1016}
\showcasevscodeiconfiletype{vue}{1017}
\showcasevscodeiconfiletype{vueconfig}{1018}
\showcasevscodeiconfiletype{vyper}{1019}
\showcasevscodeiconfiletype{wallaby}{1020}
\showcasevscodeiconfiletype{wally}{1021}
\showcasevscodeiconfiletype{wasm}{1022}
\showcasevscodeiconfiletype{watchmanconfig}{1023}
\showcasevscodeiconfiletype{wdio}{1024}
\showcasevscodeiconfiletype{weblate}{1025}
\showcasevscodeiconfiletype{webpack}{1026}
\showcasevscodeiconfiletype{wenyan}{1027}
\showcasevscodeiconfiletype{wercker}{1028}
\showcasevscodeiconfiletype{wgsl}{1029}
\showcasevscodeiconfiletype{wikitext}{1030}
\showcasevscodeiconfiletype{windi}{1031}
\showcasevscodeiconfiletype{wit}{1032}
\showcasevscodeiconfiletype{wolfram}{1033}
\showcasevscodeiconfiletype{word}{1034}
\showcasevscodeiconfiletype{word2}{1035}
\showcasevscodeiconfiletype{wpml}{1036}
\showcasevscodeiconfiletype{wurst}{1037}
\showcasevscodeiconfiletype{wxml}{1038}
\showcasevscodeiconfiletype{wxss}{1039}
\showcasevscodeiconfiletype{wxt}{1040}
\showcasevscodeiconfiletype{xcode}{1041}
\showcasevscodeiconfiletype{xfl}{1042}
\showcasevscodeiconfiletype{xib}{1043}
\showcasevscodeiconfiletype{xliff}{1044}
\showcasevscodeiconfiletype{xmake}{1045}
\showcasevscodeiconfiletype{xml}{1046}
\showcasevscodeiconfiletype{xo}{1047}
\showcasevscodeiconfiletype{xorg}{1048}
\showcasevscodeiconfiletype{xquery}{1049}
\showcasevscodeiconfiletype{xsl}{1050}
\showcasevscodeiconfiletype{yacc}{1051}
\showcasevscodeiconfiletype{yaml}{1052}
\showcasevscodeiconfiletype{yamllint}{1053}
\showcasevscodeiconfiletype{yaml-official}{1054}
\showcasevscodeiconfiletype{yandex}{1055}
\showcasevscodeiconfiletype{yang}{1056}
\showcasevscodeiconfiletype{yarn}{1057}
\showcasevscodeiconfiletype{yeoman}{1058}
\showcasevscodeiconfiletype{zig}{1059}
\showcasevscodeiconfiletype{zip}{1060}
\showcasevscodeiconfiletype{zip2}{1061}
\end{vscodeiconscase}

\subsection{FolderType version (362 icons)}

The \textsf{PDF} file is \texttt{vscodeicons-foldertype-all.pdf}

\begin{vscodeiconscase}
\showcasevscodeiconfoldertype{android}{1}
\showcasevscodeiconfoldertype{android-opened}{2}
\showcasevscodeiconfoldertype{angular}{3}
\showcasevscodeiconfoldertype{angular-opened}{4}
\showcasevscodeiconfoldertype{api}{5}
\showcasevscodeiconfoldertype{api-opened}{6}
\showcasevscodeiconfoldertype{app}{7}
\showcasevscodeiconfoldertype{app-opened}{8}
\showcasevscodeiconfoldertype{arangodb}{9}
\showcasevscodeiconfoldertype{arangodb-opened}{10}
\showcasevscodeiconfoldertype{asset}{11}
\showcasevscodeiconfoldertype{asset-opened}{12}
\showcasevscodeiconfoldertype{astro}{13}
\showcasevscodeiconfoldertype{astro-opened}{14}
\showcasevscodeiconfoldertype{audio}{15}
\showcasevscodeiconfoldertype{audio-opened}{16}
\showcasevscodeiconfoldertype{aurelia}{17}
\showcasevscodeiconfoldertype{aurelia-opened}{18}
\showcasevscodeiconfoldertype{aws}{19}
\showcasevscodeiconfoldertype{aws-opened}{20}
\showcasevscodeiconfoldertype{azure}{21}
\showcasevscodeiconfoldertype{azurepipelines}{22}
\showcasevscodeiconfoldertype{azurepipelines-opened}{23}
\showcasevscodeiconfoldertype{azure-opened}{24}
\showcasevscodeiconfoldertype{binary}{25}
\showcasevscodeiconfoldertype{binary-opened}{26}
\showcasevscodeiconfoldertype{bloc}{27}
\showcasevscodeiconfoldertype{bloc-opened}{28}
\showcasevscodeiconfoldertype{blueprint}{29}
\showcasevscodeiconfoldertype{blueprint-opened}{30}
\showcasevscodeiconfoldertype{bot}{31}
\showcasevscodeiconfoldertype{bot-opened}{32}
\showcasevscodeiconfoldertype{bower}{33}
\showcasevscodeiconfoldertype{bower-opened}{34}
\showcasevscodeiconfoldertype{buildkite}{35}
\showcasevscodeiconfoldertype{buildkite-opened}{36}
\showcasevscodeiconfoldertype{cake}{37}
\showcasevscodeiconfoldertype{cake-opened}{38}
\showcasevscodeiconfoldertype{cargo}{39}
\showcasevscodeiconfoldertype{cargo-opened}{40}
\showcasevscodeiconfoldertype{certificate}{41}
\showcasevscodeiconfoldertype{certificate-opened}{42}
\showcasevscodeiconfoldertype{changesets}{43}
\showcasevscodeiconfoldertype{changesets-opened}{44}
\showcasevscodeiconfoldertype{chef}{45}
\showcasevscodeiconfoldertype{chef-opened}{46}
\showcasevscodeiconfoldertype{circleci}{47}
\showcasevscodeiconfoldertype{circleci-opened}{48}
\showcasevscodeiconfoldertype{claude}{49}
\showcasevscodeiconfoldertype{claude-opened}{50}
\showcasevscodeiconfoldertype{cli}{51}
\showcasevscodeiconfoldertype{client}{52}
\showcasevscodeiconfoldertype{client-opened}{53}
\showcasevscodeiconfoldertype{cli-opened}{54}
\showcasevscodeiconfoldertype{cloudflare}{55}
\showcasevscodeiconfoldertype{cloudflare-opened}{56}
\showcasevscodeiconfoldertype{cmake}{57}
\showcasevscodeiconfoldertype{cmake-opened}{58}
\showcasevscodeiconfoldertype{common}{59}
\showcasevscodeiconfoldertype{common-opened}{60}
\showcasevscodeiconfoldertype{component}{61}
\showcasevscodeiconfoldertype{component-opened}{62}
\showcasevscodeiconfoldertype{composer}{63}
\showcasevscodeiconfoldertype{composer-opened}{64}
\showcasevscodeiconfoldertype{config}{65}
\showcasevscodeiconfoldertype{config-opened}{66}
\showcasevscodeiconfoldertype{controller}{67}
\showcasevscodeiconfoldertype{controller-opened}{68}
\showcasevscodeiconfoldertype{coverage}{69}
\showcasevscodeiconfoldertype{coverage-opened}{70}
\showcasevscodeiconfoldertype{css}{71}
\showcasevscodeiconfoldertype{css2}{72}
\showcasevscodeiconfoldertype{css2-opened}{73}
\showcasevscodeiconfoldertype{css-opened}{74}
\showcasevscodeiconfoldertype{cubit}{75}
\showcasevscodeiconfoldertype{cubit-opened}{76}
\showcasevscodeiconfoldertype{cypress}{77}
\showcasevscodeiconfoldertype{cypress-opened}{78}
\showcasevscodeiconfoldertype{dapr}{79}
\showcasevscodeiconfoldertype{dapr-opened}{80}
\showcasevscodeiconfoldertype{dart}{81}
\showcasevscodeiconfoldertype{dart-opened}{82}
\showcasevscodeiconfoldertype{datadog}{83}
\showcasevscodeiconfoldertype{datadog-opened}{84}
\showcasevscodeiconfoldertype{db}{85}
\showcasevscodeiconfoldertype{db-opened}{86}
\showcasevscodeiconfoldertype{debian}{87}
\showcasevscodeiconfoldertype{debian-opened}{88}
\showcasevscodeiconfoldertype{dependabot}{89}
\showcasevscodeiconfoldertype{dependabot-opened}{90}
\showcasevscodeiconfoldertype{devcontainer}{91}
\showcasevscodeiconfoldertype{devcontainer-opened}{92}
\showcasevscodeiconfoldertype{dist}{93}
\showcasevscodeiconfoldertype{dist-opened}{94}
\showcasevscodeiconfoldertype{docker}{95}
\showcasevscodeiconfoldertype{docker-opened}{96}
\showcasevscodeiconfoldertype{docs}{97}
\showcasevscodeiconfoldertype{docs-opened}{98}
\showcasevscodeiconfoldertype{e2e}{99}
\showcasevscodeiconfoldertype{e2e-opened}{100}
\showcasevscodeiconfoldertype{elasticbeanstalk}{101}
\showcasevscodeiconfoldertype{elasticbeanstalk-opened}{102}
\showcasevscodeiconfoldertype{electron}{103}
\showcasevscodeiconfoldertype{electron-opened}{104}
\showcasevscodeiconfoldertype{expo}{105}
\showcasevscodeiconfoldertype{expo-opened}{106}
\showcasevscodeiconfoldertype{favicon}{107}
\showcasevscodeiconfoldertype{favicon-opened}{108}
\showcasevscodeiconfoldertype{flow}{109}
\showcasevscodeiconfoldertype{flow-opened}{110}
\showcasevscodeiconfoldertype{fonts}{111}
\showcasevscodeiconfoldertype{fonts-opened}{112}
\showcasevscodeiconfoldertype{frontcommerce}{113}
\showcasevscodeiconfoldertype{frontcommerce-opened}{114}
\showcasevscodeiconfoldertype{gcp}{115}
\showcasevscodeiconfoldertype{gcp-opened}{116}
\showcasevscodeiconfoldertype{git}{117}
\showcasevscodeiconfoldertype{github}{118}
\showcasevscodeiconfoldertype{github-opened}{119}
\showcasevscodeiconfoldertype{gitlab}{120}
\showcasevscodeiconfoldertype{gitlab-opened}{121}
\showcasevscodeiconfoldertype{git-opened}{122}
\showcasevscodeiconfoldertype{gradle}{123}
\showcasevscodeiconfoldertype{gradle-opened}{124}
\showcasevscodeiconfoldertype{graphql}{125}
\showcasevscodeiconfoldertype{graphql-opened}{126}
\showcasevscodeiconfoldertype{grunt}{127}
\showcasevscodeiconfoldertype{grunt-opened}{128}
\showcasevscodeiconfoldertype{gulp}{129}
\showcasevscodeiconfoldertype{gulp-opened}{130}
\showcasevscodeiconfoldertype{haxelib}{131}
\showcasevscodeiconfoldertype{haxelib-opened}{132}
\showcasevscodeiconfoldertype{helper}{133}
\showcasevscodeiconfoldertype{helper-opened}{134}
\showcasevscodeiconfoldertype{hook}{135}
\showcasevscodeiconfoldertype{hook-opened}{136}
\showcasevscodeiconfoldertype{husky}{137}
\showcasevscodeiconfoldertype{husky-opened}{138}
\showcasevscodeiconfoldertype{idea}{139}
\showcasevscodeiconfoldertype{idea-opened}{140}
\showcasevscodeiconfoldertype{images}{141}
\showcasevscodeiconfoldertype{images-opened}{142}
\showcasevscodeiconfoldertype{include}{143}
\showcasevscodeiconfoldertype{include-opened}{144}
\showcasevscodeiconfoldertype{interfaces}{145}
\showcasevscodeiconfoldertype{interfaces-opened}{146}
\showcasevscodeiconfoldertype{ios}{147}
\showcasevscodeiconfoldertype{ios-opened}{148}
\showcasevscodeiconfoldertype{js}{149}
\showcasevscodeiconfoldertype{json}{150}
\showcasevscodeiconfoldertype{json-official}{151}
\showcasevscodeiconfoldertype{json-official-opened}{152}
\showcasevscodeiconfoldertype{json-opened}{153}
\showcasevscodeiconfoldertype{js-opened}{154}
\showcasevscodeiconfoldertype{junie}{155}
\showcasevscodeiconfoldertype{junie-opened}{156}
\showcasevscodeiconfoldertype{kubernetes}{157}
\showcasevscodeiconfoldertype{kubernetes-opened}{158}
\showcasevscodeiconfoldertype{less}{159}
\showcasevscodeiconfoldertype{less-opened}{160}
\showcasevscodeiconfoldertype{library}{161}
\showcasevscodeiconfoldertype{library-opened}{162}
\showcasevscodeiconfoldertype{light-cypress}{163}
\showcasevscodeiconfoldertype{light-cypress-opened}{164}
\showcasevscodeiconfoldertype{light-electron}{165}
\showcasevscodeiconfoldertype{light-electron-opened}{166}
\showcasevscodeiconfoldertype{light-expo}{167}
\showcasevscodeiconfoldertype{light-expo-opened}{168}
\showcasevscodeiconfoldertype{light-fonts}{169}
\showcasevscodeiconfoldertype{light-fonts-opened}{170}
\showcasevscodeiconfoldertype{light-gradle}{171}
\showcasevscodeiconfoldertype{light-gradle-opened}{172}
\showcasevscodeiconfoldertype{light-meteor}{173}
\showcasevscodeiconfoldertype{light-meteor-opened}{174}
\showcasevscodeiconfoldertype{light-mypy}{175}
\showcasevscodeiconfoldertype{light-mypy-opened}{176}
\showcasevscodeiconfoldertype{light-mysql}{177}
\showcasevscodeiconfoldertype{light-mysql-opened}{178}
\showcasevscodeiconfoldertype{light-node}{179}
\showcasevscodeiconfoldertype{light-node-opened}{180}
\showcasevscodeiconfoldertype{light-redux}{181}
\showcasevscodeiconfoldertype{light-redux-opened}{182}
\showcasevscodeiconfoldertype{light-sass}{183}
\showcasevscodeiconfoldertype{light-sass-opened}{184}
\showcasevscodeiconfoldertype{light-zed}{185}
\showcasevscodeiconfoldertype{light-zed-opened}{186}
\showcasevscodeiconfoldertype{linux}{187}
\showcasevscodeiconfoldertype{linux-opened}{188}
\showcasevscodeiconfoldertype{locale}{189}
\showcasevscodeiconfoldertype{locale-opened}{190}
\showcasevscodeiconfoldertype{log}{191}
\showcasevscodeiconfoldertype{log-opened}{192}
\showcasevscodeiconfoldertype{macos}{193}
\showcasevscodeiconfoldertype{macos-opened}{194}
\showcasevscodeiconfoldertype{mariadb}{195}
\showcasevscodeiconfoldertype{mariadb-opened}{196}
\showcasevscodeiconfoldertype{maven}{197}
\showcasevscodeiconfoldertype{maven-opened}{198}
\showcasevscodeiconfoldertype{mediawiki}{199}
\showcasevscodeiconfoldertype{mediawiki-opened}{200}
\showcasevscodeiconfoldertype{memcached}{201}
\showcasevscodeiconfoldertype{memcached-opened}{202}
\showcasevscodeiconfoldertype{meteor}{203}
\showcasevscodeiconfoldertype{meteor-opened}{204}
\showcasevscodeiconfoldertype{middleware}{205}
\showcasevscodeiconfoldertype{middleware-opened}{206}
\showcasevscodeiconfoldertype{minecraft}{207}
\showcasevscodeiconfoldertype{minecraft-opened}{208}
\showcasevscodeiconfoldertype{minikube}{209}
\showcasevscodeiconfoldertype{minikube-opened}{210}
\showcasevscodeiconfoldertype{mjml}{211}
\showcasevscodeiconfoldertype{mjml-opened}{212}
\showcasevscodeiconfoldertype{mock}{213}
\showcasevscodeiconfoldertype{mock-opened}{214}
\showcasevscodeiconfoldertype{model}{215}
\showcasevscodeiconfoldertype{model-opened}{216}
\showcasevscodeiconfoldertype{module}{217}
\showcasevscodeiconfoldertype{module-opened}{218}
\showcasevscodeiconfoldertype{mojo}{219}
\showcasevscodeiconfoldertype{mojo-opened}{220}
\showcasevscodeiconfoldertype{mongodb}{221}
\showcasevscodeiconfoldertype{mongodb-opened}{222}
\showcasevscodeiconfoldertype{mypy}{223}
\showcasevscodeiconfoldertype{mypy-opened}{224}
\showcasevscodeiconfoldertype{mysql}{225}
\showcasevscodeiconfoldertype{mysql-opened}{226}
\showcasevscodeiconfoldertype{netlify}{227}
\showcasevscodeiconfoldertype{netlify-opened}{228}
\showcasevscodeiconfoldertype{next}{229}
\showcasevscodeiconfoldertype{next-opened}{230}
\showcasevscodeiconfoldertype{nginx}{231}
\showcasevscodeiconfoldertype{nginx-opened}{232}
\showcasevscodeiconfoldertype{nix}{233}
\showcasevscodeiconfoldertype{nix-opened}{234}
\showcasevscodeiconfoldertype{node}{235}
\showcasevscodeiconfoldertype{node-opened}{236}
\showcasevscodeiconfoldertype{notebooks}{237}
\showcasevscodeiconfoldertype{notebooks-opened}{238}
\showcasevscodeiconfoldertype{notification}{239}
\showcasevscodeiconfoldertype{notification-opened}{240}
\showcasevscodeiconfoldertype{nuget}{241}
\showcasevscodeiconfoldertype{nuget-opened}{242}
\showcasevscodeiconfoldertype{nuxt}{243}
\showcasevscodeiconfoldertype{nuxt-opened}{244}
\showcasevscodeiconfoldertype{package}{245}
\showcasevscodeiconfoldertype{package-opened}{246}
\showcasevscodeiconfoldertype{paket}{247}
\showcasevscodeiconfoldertype{paket-opened}{248}
\showcasevscodeiconfoldertype{php}{249}
\showcasevscodeiconfoldertype{php-opened}{250}
\showcasevscodeiconfoldertype{platformio}{251}
\showcasevscodeiconfoldertype{platformio-opened}{252}
\showcasevscodeiconfoldertype{plugin}{253}
\showcasevscodeiconfoldertype{plugin-opened}{254}
\showcasevscodeiconfoldertype{prisma}{255}
\showcasevscodeiconfoldertype{prisma-opened}{256}
\showcasevscodeiconfoldertype{private}{257}
\showcasevscodeiconfoldertype{private-opened}{258}
\showcasevscodeiconfoldertype{public}{259}
\showcasevscodeiconfoldertype{public-opened}{260}
\showcasevscodeiconfoldertype{pytest}{261}
\showcasevscodeiconfoldertype{pytest-opened}{262}
\showcasevscodeiconfoldertype{python}{263}
\showcasevscodeiconfoldertype{python-opened}{264}
\showcasevscodeiconfoldertype{ravendb}{265}
\showcasevscodeiconfoldertype{ravendb-opened}{266}
\showcasevscodeiconfoldertype{redis}{267}
\showcasevscodeiconfoldertype{redis-opened}{268}
\showcasevscodeiconfoldertype{redux}{269}
\showcasevscodeiconfoldertype{redux-opened}{270}
\showcasevscodeiconfoldertype{route}{271}
\showcasevscodeiconfoldertype{route-opened}{272}
\showcasevscodeiconfoldertype{sass}{273}
\showcasevscodeiconfoldertype{sass-opened}{274}
\showcasevscodeiconfoldertype{script}{275}
\showcasevscodeiconfoldertype{script-opened}{276}
\showcasevscodeiconfoldertype{seedkit}{277}
\showcasevscodeiconfoldertype{seedkit-opened}{278}
\showcasevscodeiconfoldertype{server}{279}
\showcasevscodeiconfoldertype{server-opened}{280}
\showcasevscodeiconfoldertype{services}{281}
\showcasevscodeiconfoldertype{services-opened}{282}
\showcasevscodeiconfoldertype{shared}{283}
\showcasevscodeiconfoldertype{shared-opened}{284}
\showcasevscodeiconfoldertype{snaplet}{285}
\showcasevscodeiconfoldertype{snaplet-opened}{286}
\showcasevscodeiconfoldertype{spin}{287}
\showcasevscodeiconfoldertype{spin-opened}{288}
\showcasevscodeiconfoldertype{src}{289}
\showcasevscodeiconfoldertype{src-opened}{290}
\showcasevscodeiconfoldertype{sso}{291}
\showcasevscodeiconfoldertype{sso-opened}{292}
\showcasevscodeiconfoldertype{story}{293}
\showcasevscodeiconfoldertype{story-opened}{294}
\showcasevscodeiconfoldertype{style}{295}
\showcasevscodeiconfoldertype{style-opened}{296}
\showcasevscodeiconfoldertype{supabase}{297}
\showcasevscodeiconfoldertype{supabase-opened}{298}
\showcasevscodeiconfoldertype{svelte}{299}
\showcasevscodeiconfoldertype{svelte-opened}{300}
\showcasevscodeiconfoldertype{tauri}{301}
\showcasevscodeiconfoldertype{tauri-opened}{302}
\showcasevscodeiconfoldertype{temp}{303}
\showcasevscodeiconfoldertype{template}{304}
\showcasevscodeiconfoldertype{template-opened}{305}
\showcasevscodeiconfoldertype{temp-opened}{306}
\showcasevscodeiconfoldertype{test}{307}
\showcasevscodeiconfoldertype{test-opened}{308}
\showcasevscodeiconfoldertype{theme}{309}
\showcasevscodeiconfoldertype{theme-opened}{310}
\showcasevscodeiconfoldertype{tools}{311}
\showcasevscodeiconfoldertype{tools-opened}{312}
\showcasevscodeiconfoldertype{travis}{313}
\showcasevscodeiconfoldertype{travis-opened}{314}
\showcasevscodeiconfoldertype{trunk}{315}
\showcasevscodeiconfoldertype{trunk-opened}{316}
\showcasevscodeiconfoldertype{turbo}{317}
\showcasevscodeiconfoldertype{turbo-opened}{318}
\showcasevscodeiconfoldertype{typescript}{319}
\showcasevscodeiconfoldertype{typescript-opened}{320}
\showcasevscodeiconfoldertype{typings}{321}
\showcasevscodeiconfoldertype{typings2}{322}
\showcasevscodeiconfoldertype{typings2-opened}{323}
\showcasevscodeiconfoldertype{typings-opened}{324}
\showcasevscodeiconfoldertype{vagrant}{325}
\showcasevscodeiconfoldertype{vagrant-opened}{326}
\showcasevscodeiconfoldertype{vercel}{327}
\showcasevscodeiconfoldertype{vercel-opened}{328}
\showcasevscodeiconfoldertype{video}{329}
\showcasevscodeiconfoldertype{video-opened}{330}
\showcasevscodeiconfoldertype{view}{331}
\showcasevscodeiconfoldertype{view-opened}{332}
\showcasevscodeiconfoldertype{vs}{333}
\showcasevscodeiconfoldertype{vs2}{334}
\showcasevscodeiconfoldertype{vs2-opened}{335}
\showcasevscodeiconfoldertype{vscode}{336}
\showcasevscodeiconfoldertype{vscode2}{337}
\showcasevscodeiconfoldertype{vscode2-opened}{338}
\showcasevscodeiconfoldertype{vscode3}{339}
\showcasevscodeiconfoldertype{vscode3-opened}{340}
\showcasevscodeiconfoldertype{vscode-opened}{341}
\showcasevscodeiconfoldertype{vscode-test}{342}
\showcasevscodeiconfoldertype{vscode-test2}{343}
\showcasevscodeiconfoldertype{vscode-test2-opened}{344}
\showcasevscodeiconfoldertype{vscode-test3}{345}
\showcasevscodeiconfoldertype{vscode-test3-opened}{346}
\showcasevscodeiconfoldertype{vscode-test-opened}{347}
\showcasevscodeiconfoldertype{vs-opened}{348}
\showcasevscodeiconfoldertype{webpack}{349}
\showcasevscodeiconfoldertype{webpack-opened}{350}
\showcasevscodeiconfoldertype{windows}{351}
\showcasevscodeiconfoldertype{windows-opened}{352}
\showcasevscodeiconfoldertype{windsurf}{353}
\showcasevscodeiconfoldertype{windsurf-opened}{354}
\showcasevscodeiconfoldertype{wit}{355}
\showcasevscodeiconfoldertype{wit-opened}{356}
\showcasevscodeiconfoldertype{www}{357}
\showcasevscodeiconfoldertype{www-opened}{358}
\showcasevscodeiconfoldertype{yarn}{359}
\showcasevscodeiconfoldertype{yarn-opened}{360}
\showcasevscodeiconfoldertype{zed}{361}
\showcasevscodeiconfoldertype{zed-opened}{362}
\end{vscodeiconscase}

\end{document}